%% kapa: exam=1 round=23 year=2012 period=2 form=A subject=appraisal_law subject_ko=부동산관계법규 seq=9 number=49
%% source: Q-Net 감정평가사 기출문제 (cst00309) artlSeq=5052580 file=raw/1008/5052580/2012년 제23회 감정평가사 제1차시험 기출문제(2교시 A형).pdf
%% conversion: kordoc 4.24.51 -> tools/gen1.py (auto); review=proofread-sonnet; flags=none
\begin{question}{49}
국토의 계획 및 이용에 관한 법령상 개발행위허가에 관한 설명으로 옳은 것은?
\begin{choices}
  \item 개발행위허가를 할 때에 시장 또는 군수가「광업법」에 따른 채굴계획의 인가에 관하여 미리 관계 행정기관의 장과 협의한 사항에 대하여는 그 인가를 받은 것으로 본다.
  \item 시장 또는 군수는 지구단위계획구역으로 지정된 지역에 대하여는 3년 이내의 기간 동안 개발행위허가를 제한할 수 있으며, 개발행위허가의 제한은 그 기간을 연장할 수 없다.
  \item 해당 지방자치단체의 조례로 정하는 공공단체가 시행하는 개발행위에 대하여 시장 또는 군수는 이행보증금을 예치하게 할 수 있다.
  \item 지구단위계획을 수립한 지역에서 행하는 개발행위를 허가하려면 지방도시계획위원회의 심의를 거쳐야 한다.
  \item 개발행위허가를 받은 행정청이 새로 공공시설을 설치한 경우 새로 설치된 공공시설은 그 시설을 관리할 관리청에 유상으로 귀속된다.
\end{choices}
\end{question}
\answer{1}
